Este trabalho apresenta o desenvolvimento de uma ferramenta web de apoio à análise estrutural e ao pré-dimensionamento de pontes treliçadas construídas com palitos de picolé, com aplicação voltada à Competição de Pontes com Palitos de Picolé (CP3) da Universidade Federal Rural do Semi-Árido. A solução interpreta arquivos do FTOOL, reconstrói a geometria e a conectividade dos elementos e utiliza a biblioteca anaStruct para obter os resultados da análise estrutural. Foram implementadas rotinas de verificação de barras submetidas à tração e à compressão, considerando critérios de resistência, estabilidade e composição construtiva, para estimar quantidades de camadas e de palitos comerciais. Os módulos foram integrados a uma interface desenvolvida com Next.js, React e TypeScript, com execução das rotinas Python no navegador por meio do Pyodide. Na avaliação de um modelo de uma face de ponte, composto por 10 nós e 19 barras, os esforços axiais apresentaram diferenças de até 0,001 N em relação aos valores exibidos pelo FTOOL. O pré-dimensionamento estimou 154 camadas somadas entre as barras e 430 palitos para a face analisada. Os resultados demonstram a integração das etapas computacionais, embora as estimativas permaneçam condicionadas às hipóteses de modelagem e às propriedades adotadas para o material.

\palavraschave{pontes de palitos; análise estrutural; pré-dimensionamento; FTOOL; aplicação web.}
